\chapter{Extraction of the Impurity Operator}
\index{impurity operator!extraction|textbf}

\label{ch:impurity-operator-extraction}

\chapterstatus{proposed aligned-extraction contract; no phosphorus or boron
impurity operator is claimed.}

\section{Purpose and terminology}
\index{alignment!impurity extraction}
\index{operator difference!impurity}
\index{defect perturbation}

This chapter defines the method by which a pristine host Hamiltonian and a
doped Hamiltonian may be compared. The central rule is that an impurity
operator is not obtained by subtracting two arrays merely because they have the
same shape. The two operators must first be represented on identified finite
spaces, placed on a common scalar energy reference, and related by an explicit
coordinate map. Chapter~\ref{ch:mathematical-foundations} gives the general
operator-comparison framework; this chapter specializes it to substitutional
dopants.

\begin{definition}[Host, doped system, defect, and dopant]
\index{host system|textbf}
\index{doped system|textbf}
\index{defect|textbf}
\index{dopant|textbf}

The \emph{host} or \emph{pristine} system is the declared periodic reference
material without the substitution. The \emph{doped system} is the matched
periodic supercell in which one host atom is replaced by a dopant species. A
\emph{defect} is the resulting change relative to the host under the declared
structural, charge, spin, and relativistic branch. A \emph{dopant} is the
chemical species producing that substitution. In this monograph the dopant
label is
\begin{equation}
    d\in\{\mathrm P,\mathrm B\},
    \label{eq:impurity-dopant-labels}
\end{equation}
\begin{equationterms}{eq:impurity-dopant-labels}
  \equationterm{$d$}{the chemical dopant label.}
  \equationterm{$\mathrm P$}{substitutional phosphorus in silicon.}
  \equationterm{$\mathrm B$}{substitutional boron in silicon.}
\end{equationterms}
where $\mathrm P$ denotes substitutional phosphorus in silicon and $\mathrm B$
denotes substitutional boron in silicon.
\end{definition}

\begin{definition}[Perturbation and impurity operator]
\index{perturbation!operator|textbf}
\index{impurity operator!definition|textbf}

A \emph{perturbation operator} is a signed operator change relative to a
declared parent operator. An \emph{impurity operator} is the perturbation
obtained from pristine and doped operators only after their state spaces,
coordinates, geometry, spin structure, units, and scalar energy references have
been made compatible. It is therefore an aligned represented difference, not a
synonym for the raw dopant Hamiltonian and not automatically a scalar
potential.
\end{definition}

\begin{definition}[Onsite potential and finite-extent operator perturbation]
\index{onsite potential|textbf}
\index{finite-extent perturbation|textbf}

In an ordered localized basis, an \emph{onsite scalar potential} is diagonal in
site and internal basis labels. A \emph{finite-extent operator perturbation} may
also contain orbital-dependent onsite blocks or off-diagonal bond terms that
change hopping. Thus every onsite scalar potential is an operator
perturbation, but a perturbation containing bond, orbital-mixing, or spin-mixing
terms is not described as a scalar potential.
\end{definition}

Here and below, \emph{alignment} means the construction of an explicit map
between two represented state spaces. \emph{Gauge} means an admissible unitary
change of coordinates within one retained space. \emph{Energy-reference
alignment} means removal of a scalar offset so that energy-valued operators use
the same zero. \emph{Compatibility} means that all prerequisites for the
intended operation have been checked. These operations are distinct even when
a practical workflow performs them consecutively.

\section{Matched periodic parent problems}
\index{supercell!matched pristine and doped}
\index{periodic image}

Let the supercell be specified by positive integers
$N_1,N_2,N_3$ and primitive lattice vectors
$\mathbf a_1,\mathbf a_2,\mathbf a_3$. Its lattice vectors are
\begin{equation}
    \mathbf L_{\mu}=N_{\mu}\mathbf a_{\mu},
    \qquad \mu\in\{1,2,3\}.
    \label{eq:impurity-supercell-vectors}
\end{equation}
\begin{equationterms}{eq:impurity-supercell-vectors}
  \equationterm{$\mathbf L_{\mu}$}{the supercell lattice vector in direction $\mu$.}
  \equationterm{$N_{\mu}$}{the positive number of primitive repeats in direction $\mu$.}
  \equationterm{$\mathbf a_{\mu}$}{the primitive lattice vector in direction $\mu$.}
  \equationterm{$\mu$}{a lattice-direction index in $\{1,2,3\}$.}
\end{equationterms}
Here $\mu$ labels a lattice direction, $N_{\mu}$ is the number of primitive
repeats in that direction, and $\mathbf L_{\mu}$ is the corresponding
supercell vector. The collective label $\mathbf L$ denotes the supercell shape
and lattice vectors, not a single vector.

At supercell crystal momentum $\mathbf K$, let
\begin{equation}
    \widehat H_{b,\mathbf L}(\mathbf K)
    \quad\text{and}\quad
    \widehat H_{d,\mathbf L}(\mathbf K)
    \label{eq:impurity-parent-operators}
\end{equation}
\begin{equationterms}{eq:impurity-parent-operators}
  \equationterm{$\widehat H_{b,\mathbf L}(\mathbf K)$}{the abstract pristine Hamiltonian at supercell momentum $\mathbf K$.}
  \equationterm{$\widehat H_{d,\mathbf L}(\mathbf K)$}{the abstract doped Hamiltonian at the same momentum.}
  \equationterm{$b$}{the pristine bulk branch; $d$ is the dopant branch.}
  \equationterm{$\mathbf L$}{the complete supercell geometry and $\mathbf K$ its reduced-zone crystal momentum.}
\end{equationterms}
denote the pristine and doped Hamiltonian operators. The subscript $b$ means
bulk or pristine, $d$ is the dopant label, $\mathbf L$ identifies the
supercell, and $\mathbf K$ belongs to the reduced supercell Brillouin zone. The
operators must belong to matched physical branches: exchange--correlation
approximation, host pseudopotential, lattice vectors, boundary conditions,
relativistic treatment, and applicable spin convention must agree. A doped
supercell remains a periodic array of dopants at finite $\mathbf L$; it is not
an isolated impurity.

The doped internal coordinates may differ from the pristine coordinates when
local atomic relaxation is included. That relaxation is part of the selected
physical defect branch, but it makes the site correspondence nontrivial. A
\emph{site map} is the declared correspondence between pristine sites and the
host-derived sites of the doped supercell. The substituted site and every
unmatched or symmetry-related site must be identified rather than inferred
from matrix position.

\section{Retained spaces and finite representations}
\index{retained subspace!impurity extraction}
\index{represented operator!impurity extraction}

Let
\begin{equation}
    \mathcal S_{b,\mathbf L}(\mathbf K)
    \subseteq \mathcal H_{b,\mathbf L}(\mathbf K),
    \qquad
    \mathcal S_{d,\mathbf L}(\mathbf K)
    \subseteq \mathcal H_{d,\mathbf L}(\mathbf K)
    \label{eq:impurity-retained-space-inclusions}
\end{equation}
\begin{equationterms}{eq:impurity-retained-space-inclusions}
  \equationterm{$\mathcal S_{b,\mathbf L}(\mathbf K)$}{the finite pristine retained subspace.}
  \equationterm{$\mathcal S_{d,\mathbf L}(\mathbf K)$}{the finite doped retained subspace.}
  \equationterm{$\mathcal H_{b,\mathbf L}(\mathbf K)$}{the ambient pristine one-particle space.}
  \equationterm{$\mathcal H_{d,\mathbf L}(\mathbf K)$}{the ambient doped one-particle space.}
  \equationterm{$\subseteq$}{inclusion of a retained subspace in its ambient space.}
\end{equationterms}
be the pristine and doped retained subspaces. The symbols $\mathcal H_b$ and
$\mathcal H_d$ denote the ambient one-particle spaces, while $\mathcal S_b$ and
$\mathcal S_d$ denote the finite subspaces selected for comparison. For direct
unitary alignment, their dimensions must agree:
\begin{equation}
    \dim\mathcal S_{b,\mathbf L}(\mathbf K)
    =
    \dim\mathcal S_{d,\mathbf L}(\mathbf K)
    =M.
    \label{eq:impurity-equal-retained-dimensions}
\end{equation}
\begin{equationterms}{eq:impurity-equal-retained-dimensions}
  \equationterm{$\dim$}{the dimension of a finite retained space.}
  \equationterm{$\mathcal S_{b,\mathbf L}(\mathbf K)$ and $\mathcal S_{d,\mathbf L}(\mathbf K)$}{the pristine and doped retained spaces.}
  \equationterm{$M$}{their common finite represented dimension.}
\end{equationterms}
The integer $M$ is the represented dimension of one supercell fiber. Equal $M$
is necessary but not sufficient for compatibility.

Choose ordered orthonormal bases
\begin{equation}
    \mathcal B_b=\{\lvert b_j\rangle\}_{j=1}^{M},
    \qquad
    \mathcal B_d=\{\lvert d_i\rangle\}_{i=1}^{M}.
    \label{eq:impurity-ordered-bases}
\end{equation}
\begin{equationterms}{eq:impurity-ordered-bases}
  \equationterm{$\mathcal B_b$ and $\mathcal B_d$}{the ordered pristine and doped orthonormal bases.}
  \equationterm{$\lvert b_j\rangle$ and $\lvert d_i\rangle$}{their pristine and doped basis vectors.}
  \equationterm{$i,j$}{basis-coordinate indices running from $1$ through $M$.}
  \equationterm{$M$}{the common number of basis vectors.}
\end{equationterms}
The vectors $\lvert b_j\rangle$ and $\lvert d_i\rangle$ are basis states in the
pristine and doped retained spaces; $i$ and $j$ are coordinate indices. The
represented matrices are
\begin{align}
    [H_b]_{ij}
    &=
    \langle b_i\lvert
    \widehat H_{b,\mathbf L}(\mathbf K)
    \rvert b_j\rangle,
    \\
    [H_d]_{ij}
    &=
    \langle d_i\lvert
    \widehat H_{d,\mathbf L}(\mathbf K)
    \rvert d_j\rangle.
    \label{eq:impurity-represented-parent-matrices}
\end{align}
\begin{equationterms}{eq:impurity-represented-parent-matrices}
  \equationterm{$(H_b)_{ij}$ and $(H_d)_{ij}$}{matrix entries of the pristine and doped represented Hamiltonians.}
  \equationterm{$\widehat H_{b,\mathbf L}(\mathbf K)$ and $\widehat H_{d,\mathbf L}(\mathbf K)$}{the corresponding abstract operators.}
  \equationterm{$\lvert b_i\rangle,\lvert b_j\rangle,\lvert d_i\rangle,\lvert d_j\rangle$}{ordered basis vectors used to take matrix elements.}
  \equationterm{$\langle\cdot\lvert\cdot\rvert\cdot\rangle$}{the Hilbert-space matrix-element pairing.}
  \equationterm{$i,j$}{row and column indices; $b,d,\mathbf L,\mathbf K$ retain the meanings defined above.}
\end{equationterms}
Thus $H_b,H_d\in\mathbb C^{M\times M}$ are finite matrix representations, not
the abstract operators themselves. Their interpretation additionally requires
the basis identities, site and orbital ordering, geometry, Fourier convention,
energy unit, scalar energy zero, spin representation, and gauge.

\section{Compatibility gate}
\index{compatibility!impurity extraction}
\index{stop condition}

Before alignment or subtraction, the comparison must establish all of the
following:
\begin{enumerate}
    \item matched physical parent branches and boundary conditions;
    \item matched supercell shape and a declared geometry and site map;
    \item compatible retained spaces and represented dimensions;
    \item declared orbital, spin, and basis ordering;
    \item common energy units and Fourier conventions;
    \item a justified scalar energy-reference procedure; and
    \item an explicit alignment direction and admissible gauge class.
\end{enumerate}
A failed item is a stop condition. The method does not repair incompatible
spaces by padding matrices, discarding unexplained states, sorting eigenvalues,
or fitting an arbitrary unitary solely to reduce the final difference.

When both retained frames are available in one ambient coordinate system, let
$B_b,B_d\in\mathbb C^{N_{\mathrm{amb}}\times M}$ contain their orthonormal
basis vectors as columns. The integer $N_{\mathrm{amb}}$ is the dimension of
the common ambient discretization. The overlap matrix is
\begin{equation}
    C=B_d^{\dagger}B_b\in\mathbb C^{M\times M}.
    \label{eq:impurity-frame-overlap}
\end{equation}
\begin{equationterms}{eq:impurity-frame-overlap}
  \equationterm{$C$}{the overlap matrix from pristine coordinates to doped-frame overlaps.}
  \equationterm{$B_b$ and $B_d$}{matrices whose columns are the pristine and doped orthonormal frame vectors.}
  \equationterm{$\dagger$}{the conjugate transpose.}
  \equationterm{$\mathbb C^{M\times M}$}{the space of complex $M$-by-$M$ matrices.}
  \equationterm{$M$}{the common retained dimension.}
\end{equationterms}
The matrix $C$ measures overlap between doped and pristine retained frames; it
is not a Hamiltonian. If
\begin{equation}
    C=W\Sigma V^{\dagger}
    \label{eq:impurity-overlap-svd}
\end{equation}
\begin{equationterms}{eq:impurity-overlap-svd}
  \equationterm{$C$}{the retained-frame overlap matrix.}
  \equationterm{$W$ and $V$}{unitary left and right singular-vector matrices.}
  \equationterm{$\Sigma$}{the diagonal matrix of nonnegative singular values.}
  \equationterm{$\dagger$}{the conjugate transpose.}
\end{equationterms}
is a singular-value decomposition, then $W$ and $V$ are unitary matrices and
$\Sigma=\operatorname{diag}(\sigma_1,\ldots,\sigma_M)$ contains singular values
$0\leq\sigma_j\leq1$. The principal angles are
\begin{equation}
    \theta_j=\arccos\sigma_j,
    \qquad j\in\{1,\ldots,M\}.
    \label{eq:impurity-principal-angles}
\end{equation}
\begin{equationterms}{eq:impurity-principal-angles}
  \equationterm{$\theta_j$}{the $j$th principal angle between retained spaces.}
  \equationterm{$\sigma_j$}{the $j$th singular value of the overlap matrix.}
  \equationterm{$\arccos$}{the inverse cosine on $[0,1]$.}
  \equationterm{$j$}{an index in $\{1,\ldots,M\}$; $M$ is the retained dimension.}
\end{equationterms}
Each $\theta_j$ diagnoses one direction of subspace agreement. A singular value
near zero signals an unresolved or nearly orthogonal direction; it cannot be
fixed by relabelling the basis.

For compatible full-rank frames, the unitary polar factor
\begin{equation}
    U_{\mathrm{pol}}=WV^{\dagger}
    \label{eq:impurity-polar-alignment}
\end{equation}
\begin{equationterms}{eq:impurity-polar-alignment}
  \equationterm{$U_{\mathrm{pol}}$}{the unitary polar factor mapping pristine coordinates to doped coordinates under the overlap convention.}
  \equationterm{$W$ and $V$}{the singular-vector matrices of $C$.}
  \equationterm{$\dagger$}{the conjugate transpose.}
\end{equationterms}
is the nearest unitary factor associated with $C$ in the standard
orthogonal-Procrustes construction. Here $U_{\mathrm{pol}}$ maps pristine
coordinate vectors to doped coordinate vectors under the stated overlap
convention. It is one possible alignment ingredient, not a universal physical
answer. A project alignment may additionally constrain site permutations,
orbital anchors, spin frames, crystal symmetries, smoothness in $\mathbf K$,
and periodicity across the reciprocal mesh. The selected map, its direction,
conditioning, singular values, residuals, and non-uniqueness must be retained.
No universal singular-value or principal-angle cutoff is asserted here.

\section{Scalar energy-reference alignment}
\index{energy reference!impurity extraction}

Let $E_{0,b}$ and $E_{0,d}$ be the scalar energy zeros assigned to the pristine
and doped representations. Define
\begin{equation}
    \overline H_b=H_b-E_{0,b}I_M,
    \qquad
    \overline H_d=H_d-E_{0,d}I_M,
    \label{eq:impurity-energy-aligned-matrices}
\end{equation}
\begin{equationterms}{eq:impurity-energy-aligned-matrices}
  \equationterm{$\overline H_b$ and $\overline H_d$}{pristine and doped matrices after scalar energy alignment.}
  \equationterm{$H_b$ and $H_d$}{the source represented Hamiltonian matrices.}
  \equationterm{$E_{0,b}$ and $E_{0,d}$}{their declared scalar energy zeros.}
  \equationterm{$I_M$}{the identity matrix on the $M$-dimensional represented space.}
  \equationterm{$b$ and $d$}{pristine and doped branches; the overline denotes energy alignment.}
\end{equationterms}
where $I_M$ is the $M\times M$ identity matrix and the overline denotes an
energy-aligned representation. Both $E_{0,b}$ and $E_{0,d}$ have the same
energy unit as the matrix entries. Only their difference affects the extracted
operator, but both source conventions must be recorded.

One general diagnostic for a scalar offset uses a declared reference-sector
projector $P_{\mathrm{ref}}$. Suppose a provisional coordinate alignment $U$
has been supplied and define
\begin{equation}
    X=H_d-UH_bU^{\dagger}.
    \label{eq:impurity-provisional-difference}
\end{equation}
\begin{equationterms}{eq:impurity-provisional-difference}
  \equationterm{$X$}{the provisional uncorrected difference in doped coordinates.}
  \equationterm{$H_d$ and $H_b$}{doped and pristine represented Hamiltonians.}
  \equationterm{$U$}{a supplied provisional map from pristine to doped coordinates.}
  \equationterm{$\dagger$}{the conjugate transpose.}
\end{equationterms}
The matrix $X$ is the uncorrected doped-space difference. If
$P_{\mathrm{ref}}=P_{\mathrm{ref}}^{\dagger}=P_{\mathrm{ref}}^2$ projects onto
a sector assumed not to contain the localized defect signal, the least-squares
scalar offset on that sector is
\begin{equation}
    c_{\mathrm{ref}}
    =
    \frac{\operatorname{Re}\operatorname{Tr}
    (P_{\mathrm{ref}}XP_{\mathrm{ref}})}
    {\operatorname{rank}(P_{\mathrm{ref}})}.
    \label{eq:impurity-reference-offset}
\end{equation}
\begin{equationterms}{eq:impurity-reference-offset}
  \equationterm{$c_{\mathrm{ref}}$}{the least-squares scalar energy offset estimated on the reference sector.}
  \equationterm{$P_{\mathrm{ref}}$}{the orthogonal projector onto that declared sector.}
  \equationterm{$X$}{the provisional uncorrected doped-space difference.}
  \equationterm{$\operatorname{Tr}$ and $\operatorname{Re}$}{the matrix trace and real-part operations.}
  \equationterm{$\operatorname{rank}(P_{\mathrm{ref}})$}{the number of retained reference directions.}
\end{equationterms}
Here $\operatorname{Tr}$ is the matrix trace, $\operatorname{Re}$ takes the real
part, and $\operatorname{rank}(P_{\mathrm{ref}})$ is the number of retained
reference directions. This formula explains a class of exterior-anchor
estimators; it is not yet the frozen material estimator. It is undefined when
the reference rank is zero and can be biased when the defect perturbation
extends into the reference sector. The production estimator, spatial sampling
region, convergence test, and uncertainty remain required inputs to the dopant
program.

\section{Common-space alignment and signed extraction}
\index{alignment map!pristine to doped}
\index{operator difference!direction}

Let
\begin{equation}
    U_d:
    \mathcal S_{b,\mathbf L}(\mathbf K)
    \longrightarrow
    \mathcal S_{d,\mathbf L}(\mathbf K)
    \label{eq:impurity-identification-map}
\end{equation}
\begin{equationterms}{eq:impurity-identification-map}
  \equationterm{$U_d$}{the selected coordinate identification from pristine to doped retained space for dopant $d$.}
  \equationterm{$\mathcal S_{b,\mathbf L}(\mathbf K)$}{the pristine retained space.}
  \equationterm{$\mathcal S_{d,\mathbf L}(\mathbf K)$}{the doped retained space.}
  \equationterm{$\mathbf L$ and $\mathbf K$}{the supercell specification and supercell crystal momentum.}
\end{equationterms}
be the selected identification from pristine to doped retained space. For an
equal-rank comparison, $U_d$ is required to be unitary:
\begin{equation}
    U_d^{\dagger}U_d=I_b,
    \qquad
    U_dU_d^{\dagger}=I_d,
    \label{eq:impurity-unitary-identification}
\end{equation}
\begin{equationterms}{eq:impurity-unitary-identification}
  \equationterm{$U_d$}{the equal-rank pristine-to-doped unitary identification.}
  \equationterm{$U_d^{\dagger}$}{its adjoint and inverse.}
  \equationterm{$I_b$ and $I_d$}{identity operators on pristine and doped retained spaces.}
\end{equationterms}
where $I_b$ and $I_d$ are the identity operators on the pristine and doped
retained spaces. The dagger denotes the adjoint. If only a common sector is
identified, $U_d$ may instead be a declared partial isometry; then no full-space
impurity operator is inferred on the unidentified complement.

The pristine operator pushed into doped coordinates is
\begin{equation}
    \overline H_{b\to d}
    =
    U_d\overline H_bU_d^{\dagger}.
    \label{eq:impurity-pristine-pushforward}
\end{equation}
\begin{equationterms}{eq:impurity-pristine-pushforward}
  \equationterm{$\overline H_{b\to d}$}{the energy-aligned pristine Hamiltonian represented in doped coordinates.}
  \equationterm{$\overline H_b$}{the energy-aligned pristine matrix in its original coordinates.}
  \equationterm{$U_d$ and $U_d^{\dagger}$}{the pristine-to-doped map and its adjoint.}
  \equationterm{$b\to d$}{the direction of coordinate transport.}
\end{equationterms}
The arrow $b\to d$ states the direction of coordinate transport. The signed
impurity operator represented in doped coordinates is
\begin{equation}
    \boxed{
    \Delta H_{b\to d}
    =
    \overline H_d
    -
    U_d\overline H_bU_d^{\dagger}}
    .
    \label{eq:doped-space-impurity-operator}
\end{equation}
\begin{equationterms}{eq:doped-space-impurity-operator}
  \equationterm{$\Delta H_{b\to d}$}{the signed impurity matrix in doped coordinates.}
  \equationterm{$\overline H_d$}{the energy-aligned doped Hamiltonian.}
  \equationterm{$\overline H_b$}{the energy-aligned pristine Hamiltonian.}
  \equationterm{$U_d$ and $U_d^{\dagger}$}{the pristine-to-doped alignment and its adjoint.}
  \equationterm{$b\to d$}{the common doped-coordinate comparison direction.}
\end{equationterms}
The first operand is the doped operator, so a positive onsite entry denotes an
increase relative to the aligned host convention.

Alternatively, pull the doped operator into pristine coordinates:
\begin{equation}
    \overline H_{d\to b}
    =
    U_d^{\dagger}\overline H_dU_d,
    \label{eq:impurity-doped-pullback}
\end{equation}
\begin{equationterms}{eq:impurity-doped-pullback}
  \equationterm{$\overline H_{d\to b}$}{the energy-aligned doped Hamiltonian represented in pristine coordinates.}
  \equationterm{$\overline H_d$}{the energy-aligned doped matrix in its original coordinates.}
  \equationterm{$U_d^{\dagger}$ and $U_d$}{the pullback map and its inverse orientation.}
  \equationterm{$d\to b$}{the direction of coordinate transport.}
\end{equationterms}
and define
\begin{equation}
    \boxed{
    \Delta H_{d\to b}
    =
    U_d^{\dagger}\overline H_dU_d
    -
    \overline H_b}
    .
    \label{eq:pristine-space-impurity-operator}
\end{equation}
\begin{equationterms}{eq:pristine-space-impurity-operator}
  \equationterm{$\Delta H_{d\to b}$}{the signed impurity matrix in pristine coordinates.}
  \equationterm{$\overline H_d$ and $\overline H_b$}{the aligned doped and pristine Hamiltonians.}
  \equationterm{$U_d^{\dagger}$ and $U_d$}{the maps that pull the doped operator into pristine coordinates.}
  \equationterm{$d\to b$}{the common pristine-coordinate comparison direction.}
\end{equationterms}
If $U_d$ is unitary, the two differences describe the same aligned perturbation
in different coordinates:
\begin{equation}
    \Delta H_{b\to d}
    =
    U_d\Delta H_{d\to b}U_d^{\dagger}.
    \label{eq:impurity-coordinate-equivalence}
\end{equation}
\begin{equationterms}{eq:impurity-coordinate-equivalence}
  \equationterm{$\Delta H_{b\to d}$ and $\Delta H_{d\to b}$}{the same aligned perturbation represented in doped and pristine coordinates.}
  \equationterm{$U_d$ and $U_d^{\dagger}$}{the unitary coordinate identification and its adjoint.}
\end{equationterms}
This equality explains why the alignment direction must be recorded. The two
matrices need not be entrywise equal, although unitary-invariant norms and
spectra agree. The extraction operation itself performs no site mapping, gauge
inference, basis rotation, geometry transformation, or energy-zero estimation;
those are prerequisites represented by $U_d$, the overlines, and the
compatibility record.

\section{Localized matrix interpretation}
\index{onsite block}
\index{bond term}
\index{spin mixing}

After a localized basis has been justified, write one basis index as a pair
$(i,\alpha)$, where $i$ labels a lattice site and $\alpha$ labels an internal
degree of freedom such as an orbital or spin component. The matrix elements of
the aligned impurity operator are
\begin{equation}
    [\Delta H]_{i\alpha,j\beta}
    =
    \langle i,\alpha\lvert
    \Delta\widehat H
    \rvert j,\beta\rangle.
    \label{eq:impurity-localized-matrix-elements}
\end{equation}
\begin{equationterms}{eq:impurity-localized-matrix-elements}
  \equationterm{$(\Delta H)_{i\alpha,j\beta}$}{one represented matrix element of the aligned impurity operator.}
  \equationterm{$\Delta\widehat H$}{the corresponding abstract perturbation operator.}
  \equationterm{$i,j$}{lattice-site labels and $\alpha,\beta$ internal orbital or spin labels.}
  \equationterm{$\lvert i,\alpha\rangle$ and $\lvert j,\beta\rangle$}{localized basis states.}
\end{equationterms}
The indices $j$ and $\beta$ label the ket-side site and internal state. Terms
with $i=j$ are onsite blocks. Terms with $i\neq j$ are bond or hopping changes.
Terms with $\alpha\neq\beta$ mix internal components and may describe orbital
or spin mixing, depending on the basis.

A scalar onsite perturbation has the form
\begin{equation}
    [\Delta H_{\mathrm{pot}}]_{i\alpha,j\beta}
    =
    V_i\,\delta_{ij}\delta_{\alpha\beta},
    \label{eq:impurity-onsite-potential}
\end{equation}
\begin{equationterms}{eq:impurity-onsite-potential}
  \equationterm{$\Delta H_{\mathrm{pot}}$}{a scalar onsite perturbation matrix.}
  \equationterm{$V_i$}{the energy-valued scalar onsite change at site $i$.}
  \equationterm{$i,j$}{site labels and $\alpha,\beta$ internal labels.}
  \equationterm{$\delta_{ij}$ and $\delta_{\alpha\beta}$}{Kronecker deltas enforcing equality of their indices.}
\end{equationterms}
where $V_i$ is an energy-valued onsite coefficient and $\delta_{ij}$ and
$\delta_{\alpha\beta}$ are Kronecker deltas. This operator is diagonal in both
site and internal labels and is proportional to the identity within each
onsite internal space.

A general Hermitian finite-extent bond perturbation may be written
\begin{equation}
    \Delta H_{\mathrm{bond}}
    =
    \sum_{(i\alpha,j\beta)\in\mathcal E}
    \left(
      t_{i\alpha,j\beta}
      \lvert i,\alpha\rangle\langle j,\beta\rvert
      +
      t_{i\alpha,j\beta}^{*}
      \lvert j,\beta\rangle\langle i,\alpha\rvert
    \right).
    \label{eq:impurity-bond-perturbation}
\end{equation}
\begin{equationterms}{eq:impurity-bond-perturbation}
  \equationterm{$\Delta H_{\mathrm{bond}}$}{the Hermitian bond-change operator.}
  \equationterm{$\mathcal E$}{a finite set of directed basis-state pairs $(i\alpha,j\beta)$.}
  \equationterm{$t_{i\alpha,j\beta}$}{the complex energy-valued directed bond change.}
  \equationterm{$t_{i\alpha,j\beta}^{*}$}{its complex conjugate.}
  \equationterm{$\lvert i,\alpha\rangle\langle j,\beta\rvert$}{the operator mapping the ket state $(j,\beta)$ to $(i,\alpha)$.}
\end{equationterms}
Here $\mathcal E$ is a finite set of directed basis-state pairs,
$t_{i\alpha,j\beta}$ is a complex energy-valued bond change, and the second
term is its Hermitian conjugate. Such a perturbation modifies hopping and must
not be relabelled as a scalar potential.

\section{Locality and finite-extent diagnostics}
\index{minimum-image distance}
\index{locality!partition-resolved norm}

Let site $i$ have integer supercell coordinates
$\mathbf n_i=(n_{i1},n_{i2},n_{i3})$, with
$n_{i\mu}\in\{0,\ldots,N_{\mu}-1\}$. Let $o$ denote the declared defect-origin
site. The minimum-image separation in direction $\mu$ is
\begin{equation}
    \delta_{\mu}(i,o)
    =
    \min\left(
      \lvert n_{i\mu}-n_{o\mu}\rvert,
      N_{\mu}-\lvert n_{i\mu}-n_{o\mu}\rvert
    \right).
    \label{eq:impurity-minimum-image-separation}
\end{equation}
\begin{equationterms}{eq:impurity-minimum-image-separation}
  \equationterm{$\delta_{\mu}(i,o)$}{the minimum-image site separation in lattice direction $\mu$.}
  \equationterm{$i$ and $o$}{the evaluated site and declared defect-origin site.}
  \equationterm{$n_{i\mu}$ and $n_{o\mu}$}{their integer coordinates in direction $\mu$.}
  \equationterm{$N_{\mu}$}{the periodic site count in direction $\mu$.}
  \equationterm{$\lvert\cdot\rvert$ and $\min$}{absolute value and selection of the shorter periodic displacement.}
\end{equationterms}
The minimum-image Chebyshev distance is
\begin{equation}
    d_{\infty}^{\mathrm{MI}}(i,o)
    =
    \max_{\mu}\delta_{\mu}(i,o).
    \label{eq:impurity-minimum-image-chebyshev}
\end{equation}
\begin{equationterms}{eq:impurity-minimum-image-chebyshev}
  \equationterm{$d_{\infty}^{\mathrm{MI}}(i,o)$}{the minimum-image Chebyshev distance between site $i$ and origin $o$.}
  \equationterm{$\delta_{\mu}(i,o)$}{the minimum-image separation in direction $\mu$.}
  \equationterm{$\max_{\mu}$}{the maximum over lattice directions.}
\end{equationterms}
The minimum implements periodic wrapping, while the maximum assigns a lattice
shell. Other metrics may be used only when declared; they define different
spatial partitions.

For a nonnegative integer radius $r$, define the core projector
\begin{equation}
    P_r
    =
    \sum_{i:\,d_{\infty}^{\mathrm{MI}}(i,o)\leq r}
    \sum_{\alpha}
    \lvert i,\alpha\rangle\langle i,\alpha\rvert,
    \qquad
    Q_r=I-P_r.
    \label{eq:impurity-core-projectors}
\end{equation}
\begin{equationterms}{eq:impurity-core-projectors}
  \equationterm{$P_r$}{the orthogonal projector onto all localized basis states within radius $r$ of origin $o$.}
  \equationterm{$Q_r$}{the complementary exterior projector.}
  \equationterm{$r$}{a nonnegative integer shell radius.}
  \equationterm{$i$ and $\alpha$}{a site and its internal basis-state label.}
  \equationterm{$d_{\infty}^{\mathrm{MI}}(i,o)$}{the declared minimum-image distance and $I$ the full represented identity.}
\end{equationterms}
The projector $P_r$ selects all internal basis states on sites within the core;
$Q_r$ selects its exterior. The partition-resolved Frobenius norms are
\begin{align}
    \varepsilon_{\mathrm{core}}(r)
    &=\lVert P_r\Delta H P_r\rVert_{\mathrm F},
    \\
    \varepsilon_{\mathrm{ext}}(r)
    &=\lVert Q_r\Delta H Q_r\rVert_{\mathrm F},
    \\
    \varepsilon_{\mathrm{cross}}(r)
    &=
    \left(
      \lVert P_r\Delta H Q_r\rVert_{\mathrm F}^{2}
      +
      \lVert Q_r\Delta H P_r\rVert_{\mathrm F}^{2}
    \right)^{1/2}.
    \label{eq:impurity-partitioned-locality-norms}
\end{align}
\begin{equationterms}{eq:impurity-partitioned-locality-norms}
  \equationterm{$\varepsilon_{\mathrm{core}}(r)$}{the Frobenius norm of the core--core block.}
  \equationterm{$\varepsilon_{\mathrm{ext}}(r)$}{the Frobenius norm of the exterior--exterior block.}
  \equationterm{$\varepsilon_{\mathrm{cross}}(r)$}{the combined Frobenius norm of both core--exterior blocks.}
  \equationterm{$P_r$ and $Q_r$}{the core and exterior projectors at radius $r$.}
  \equationterm{$\Delta H$}{the aligned represented impurity operator and $\lVert\cdot\rVert_{\mathrm F}$ the Frobenius norm.}
\end{equationterms}
The core norm measures perturbation wholly inside the core. The exterior norm
measures perturbation wholly outside it. The cross norm measures couplings that
cross the core boundary. All three have energy units.

For shell $s\geq0$, define $P_{-1}=0$ and
\begin{equation}
    S_s=P_s-P_{s-1},
    \qquad
    \eta_s=\lVert S_s\Delta H\rVert_{\mathrm F}.
    \label{eq:impurity-shell-norm}
\end{equation}
\begin{equationterms}{eq:impurity-shell-norm}
  \equationterm{$S_s$}{the projector onto lattice shell $s$.}
  \equationterm{$P_s$ and $P_{s-1}$}{cumulative core projectors through shells $s$ and $s-1$; $P_{-1}=0$.}
  \equationterm{$\eta_s$}{the row-resolved Frobenius norm on shell $s$.}
  \equationterm{$\Delta H$}{the aligned represented impurity operator.}
  \equationterm{$s$}{a nonnegative integer shell index.}
\end{equationterms}
The projector $S_s$ selects rows associated with sites at shell $s$, and
$\eta_s$ is the row-resolved shell norm. Because a general unitary can mix
sites, $\eta_s$ and the partitioned norms are meaningful only in the declared
localized gauge or under a restricted site-local gauge class.

A finite-extent acceptance rule at radius $r$ takes explicit nonnegative
energy-valued tolerances $\tau_{\mathrm{ext}}$ and
$\tau_{\mathrm{cross}}$ and requires
\begin{equation}
    \varepsilon_{\mathrm{ext}}(r)\leq\tau_{\mathrm{ext}},
    \qquad
    \varepsilon_{\mathrm{cross}}(r)\leq\tau_{\mathrm{cross}}.
    \label{eq:impurity-finite-extent-criterion}
\end{equation}
\begin{equationterms}{eq:impurity-finite-extent-criterion}
  \equationterm{$\varepsilon_{\mathrm{ext}}(r)$}{the exterior--exterior residual at radius $r$.}
  \equationterm{$\varepsilon_{\mathrm{cross}}(r)$}{the core--exterior residual at radius $r$.}
  \equationterm{$\tau_{\mathrm{ext}}$ and $\tau_{\mathrm{cross}}$}{authored nonnegative energy-valued acceptance tolerances.}
  \equationterm{$r$}{the declared nonnegative core radius.}
\end{equationterms}
The tolerances, radius, metric, basis, and gauge are part of the claim. Passing
this finite represented criterion does not prove compact support in a continuum
or infinite crystal.

\section{Supercell convergence and periodic-image contamination}
\index{supercell convergence!impurity operator}
\index{periodic-image contamination}

Let $\Delta H_{\mathbf L}$ and $\Delta H_{\mathbf L'}$ be aligned impurity
operators from two supercells. They cannot be subtracted until a common local
space has been identified. Let
\begin{equation}
    J_{\mathbf L\to c}:\mathcal S_{\mathbf L}\longrightarrow\mathcal S_c,
    \qquad
    J_{\mathbf L'\to c}:\mathcal S_{\mathbf L'}\longrightarrow\mathcal S_c
    \label{eq:impurity-common-space-embeddings}
\end{equation}
\begin{equationterms}{eq:impurity-common-space-embeddings}
  \equationterm{$J_{\mathbf L\to c}$ and $J_{\mathbf L'\to c}$}{declared maps from two supercell spaces into one common comparison space.}
  \equationterm{$\mathcal S_{\mathbf L}$ and $\mathcal S_{\mathbf L'}$}{represented spaces for supercells $\mathbf L$ and $\mathbf L'$.}
  \equationterm{$\mathcal S_c$}{the common local comparison space.}
  \equationterm{$\mathbf L\to c$ and $\mathbf L'\to c$}{map directions, not arithmetic operations.}
\end{equationterms}
be declared restriction or embedding maps into a common comparison space
$\mathcal S_c$. The core stability residual is
\begin{equation}
    \varepsilon_{\mathrm{size}}(r;\mathbf L,\mathbf L')
    =
    \left\lVert
      P_r
      \left(
        J_{\mathbf L\to c}\Delta H_{\mathbf L}J_{\mathbf L\to c}^{\dagger}
        -
        J_{\mathbf L'\to c}\Delta H_{\mathbf L'}J_{\mathbf L'\to c}^{\dagger}
      \right)
      P_r
    \right\rVert_{\mathrm F}.
    \label{eq:impurity-supercell-stability-residual}
\end{equation}
\begin{equationterms}{eq:impurity-supercell-stability-residual}
  \equationterm{$\varepsilon_{\mathrm{size}}(r;\mathbf L,\mathbf L')$}{the core Frobenius discrepancy between two supercell extractions.}
  \equationterm{$\Delta H_{\mathbf L}$ and $\Delta H_{\mathbf L'}$}{aligned impurity matrices extracted from supercells $\mathbf L$ and $\mathbf L'$.}
  \equationterm{$J_{\mathbf L\to c}$ and $J_{\mathbf L'\to c}$}{their maps into the common comparison space; daggers denote adjoints.}
  \equationterm{$P_r$}{the common core projector at radius $r$.}
  \equationterm{$\lVert\cdot\rVert_{\mathrm F}$}{the Frobenius norm.}
\end{equationterms}
The maps $J$ define which sites, orbitals, and spin components are compared;
$\mathcal S_c$ is their common represented space. The same radius $r$ and core
projector $P_r$ must refer to the same physical neighborhood in both
representations.

A convergence study must vary both supercell size and shape. Increasing size
changes dopant-image separation. Changing shape probes anisotropic image
coupling and exposes conclusions that depend accidentally on a square or cubic
sequence. Defect-derived eigenvalues must also be tracked over supercell
momentum $\mathbf K$: their bandwidth is evidence of periodic-image coupling
and must not be collapsed into one nominal isolated bound-state energy.

\section{Verification residuals and evidence classes}
\index{verification!impurity extraction}
\index{recovery residual}

For a synthetic control with a known planted perturbation
$\Delta H_{\star}$ and extracted result $\Delta H_{\mathrm{ext}}$, define the
recovery error
\begin{equation}
    R=\Delta H_{\mathrm{ext}}-\Delta H_{\star}.
    \label{eq:impurity-recovery-error}
\end{equation}
\begin{equationterms}{eq:impurity-recovery-error}
  \equationterm{$R$}{the extraction-recovery error matrix.}
  \equationterm{$\Delta H_{\mathrm{ext}}$}{the perturbation returned by the extraction method.}
  \equationterm{$\Delta H_{\star}$}{the independently authored planted perturbation.}
  \equationterm{$\star$}{an authored-reference marker, not complex conjugation in this equation.}
\end{equationterms}
The symbol $\star$ denotes the authored reference, not an optimized value. Two
absolute recovery measures are
\begin{equation}
    \varepsilon_{\max}=\max_{p,q}\lvert R_{pq}\rvert,
    \qquad
    \varepsilon_{\mathrm F}=\lVert R\rVert_{\mathrm F},
    \label{eq:impurity-recovery-norms}
\end{equation}
\begin{equationterms}{eq:impurity-recovery-norms}
  \equationterm{$\varepsilon_{\max}$}{the largest absolute entry of the recovery matrix.}
  \equationterm{$\varepsilon_{\mathrm F}$}{the Frobenius norm of the recovery matrix.}
  \equationterm{$R_{pq}$}{the recovery-matrix entry at represented indices $p,q$.}
  \equationterm{$p,q$}{indices over the represented basis.}
  \equationterm{$\lvert\cdot\rvert$, $\max$, and $\lVert\cdot\rVert_{\mathrm F}$}{absolute value, maximum, and Frobenius norm.}
\end{equationterms}
where $p$ and $q$ range over represented basis indices. The maximum-entry norm
locates the largest coordinate discrepancy; the Frobenius norm aggregates all
entries. These residuals verify the declared finite representation and
extraction route when compared with authored tolerances. They do not establish
that the planted model is physically adequate for silicon.

The error sources remain separate:
\begin{itemize}
    \item \emph{parent-model error} concerns the Kohn--Sham approximation,
    exchange--correlation functional, pseudopotentials, and physical branch;
    \item \emph{numerical error} concerns discretization, self-consistency,
    reciprocal meshes, Wannier convergence, and finite supercells;
    \item \emph{alignment error} concerns site, subspace, gauge, spin, and
    scalar-energy identification; and
    \item \emph{reduction error} concerns approximation of the aligned
    impurity operator by a lattice or continuum model class.
\end{itemize}
A successful software test establishes only its software contract. Numerical
verification requires an independent mathematical or computational route.
Scientific validation requires comparison with independent physical evidence
for a declared use. Sensitivity scans are not uncertainty quantification unless
uncertain inputs and their propagation have been defined.

\section{Dopant-specific interpretation}
\index{phosphorus!neutral branch}
\index{boron!spin--orbit branch}

For phosphorus, the primary branch is neutral spin-polarized
$\mathrm P_{\mathrm{Si}}^{0}$. It contains both the extra electron and the
positively charged donor core relative to silicon. Therefore
Eq.~\eqref{eq:doped-space-impurity-operator} yields the aligned perturbation of
the complete selected neutral Kohn--Sham branch. It must not be silently
relabelled as the potential of an ionized donor core.

For boron, the final acceptor branch must retain the accepted noncollinear
spinor and spin--orbit-coupled valence structure. Its pristine comparison must
use the same relativistic and spin convention. A scalar-relativistic non-SOC
boron calculation may support method development only; it is not the final
acceptor impurity operator.

For both dopants, the matrix must be extracted over a supercell sequence.
Stability of local blocks, gauge and alignment sensitivity, exterior and
cross-coupling norms, defect-band dispersion, and shape dependence are required
before an isolated-impurity interpretation is considered. One finite-cell
matrix difference is an impurity-operator candidate, not an isolated-impurity
result.

Appendix~\ref{app:impurity-model-classes} records the controlled
one-dimensional extraction, alignment, model-class, and continuum exercises.
Appendix~\ref{app:two-dimensional-defect-extraction} records the finite-periodic
two-dimensional representation, extraction, locality, symmetry, and adverse
controls. Those controlled calculations verify bounded synthetic contracts.
They do not constitute phosphorus or boron extraction, material validation,
transferability evidence, or uncertainty quantification.
